\chapter{超越单个行列式：电子关联}
\sourcepages{手写笔记第 31--35 页}
\section{平均场遗漏了什么}
HF 已考虑交换，但某个电子的位置涨落如何即时影响另一个电子，单行列式通常无法完整描述。在同一 Hamilton 算符和同一单粒子基空间内，定义关联能 $E_c=E_{\rm exact}-E_{\rm HF}\le0$。有限基 full CI 的“精确”只指该基空间内；离连续空间仍有基组误差。

拉长氢分子说明了另一种困难：远距离极限应各有一个电子留在不同原子，受限闭壳层的双占据成键轨道却含有不合理的离子型组分。此时多个行列式具有同等重要性，不能把缺失部分都当作小修正。与弱涨落对应的动态关联相比，这种近简并所需的多构型描述常称为静态关联。

\section{CI：扩大变分空间}
从 HF 参考行列式 $\Phi_0$ 出发，把占据轨道 $i$ 换成未占据轨道 $a$ 得 $\Phi_i^a$，替换两个得 $\Phi_{ij}^{ab}$。这些不是粒子沿真实时间跃迁的历史，而是多体 Hilbert 空间的基向量。展开
\begin{equation}
\Psi=c_0\Phi_0+\sum_{ia}c_i^a\Phi_i^a+\sum_{i<j,a<b}c_{ij}^{ab}\Phi_{ij}^{ab}+\cdots.
\end{equation}
用正交轨道组成的不同占据行列式互相正交。将展开代入 Schrödinger 方程，左乘 $\langle\Phi_I|$，得到 $\sum_JH_{IJ}c_J=Ec_I$。这是普通多体矩阵本征问题。

包含所有 $\binom MN$ 个行列式就是 full CI；组合数增长使其很快不可承受。截断至单、双激发得到 CISD，并保留变分上界，却一般不具尺寸一致性。若两个不相互作用片段各需要双激发，乘积波函数含有“片段 A 双激发乘片段 B 双激发”的总四激发；整体 CISD 把它丢掉，所以不能保证 $E(A+B)=E(A)+E(B)$。

\section{MP2：把剩余作用当作微扰}
取 $\hat H_0=\sum_i\hat f(i)$，$\hat V=\hat H-\hat H_0$，并暂引参数 $\lambda$。把 $E=E^{(0)}+\lambda E^{(1)}+\lambda^2E^{(2)}+\cdots$、$|\Psi\rangle=|0\rangle+\lambda|\Psi^{(1)}\rangle+\cdots$ 代入本征方程。采用 $\langle0|\Psi^{(1)}\rangle=0$，一阶方程投影到 $|0\rangle$ 给出 $E^{(1)}=\langle0|V|0\rangle$；投影到 $|K\rangle$ 给出
\begin{equation}
\langle K|\Psi^{(1)}\rangle=\frac{\langle K|V|0\rangle}{E^{(0)}_0-E^{(0)}_K}.
\end{equation}
二阶方程再投影到 $|0\rangle$，便得
\begin{equation}
E^{(2)}=\sum_{K\ne0}\frac{|\langle K|V|0\rangle|^2}{E^{(0)}_0-E^{(0)}_K}.
\end{equation}
这里分子必须连接不同态；若误写成每一项都是 $\langle0|V|0\rangle$，就丢掉了激发通道的信息。

对正则 HF 参考，轨道驻值条件使单激发矩阵元为零；一体与二体算符无法连接相差三个及以上轨道的行列式，因为至少留下一个正交轨道重叠。故只剩双激发：
\begin{equation}
E^{(2)}_{\rm MP2}=\sum_{i<j,a<b}
\frac{|\langle\Phi_{ij}^{ab}|\hat V|\Phi_0\rangle|^2}{\varepsilon_i+\varepsilon_j-\varepsilon_a-\varepsilon_b},\qquad E_{\rm MP2}=E_{\rm HF}+E^{(2)}_{\rm MP2}.
\end{equation}
分母是未扰动能量差，不是 HF 总能量与“激发行列式”相减。通常占据态低于虚态，修正为负；但能隙很小时修正可能失控。MP2 不是能量最小化，所以它可以低于 full CI 能量，并无“必在 HF 与精确解之间”的保证，也没有普适的关联能恢复百分比。

\section{耦合簇为什么使用指数}
为自动包含独立片段上同时发生的激发，写 $|\Psi\rangle=e^{\hat T}|\Phi_0\rangle$，其中 $\hat T=\hat T_1+\hat T_2+\cdots$，$\hat T_k$ 产生 $k$ 重激发。即便截断 $T=T_1+T_2$，展开指数仍包含 $T_2^2/2$ 等四激发。这正补上了截断 CI 缺少的独立乘积结构。

若不相互作用片段满足 $T=T_A+T_B$ 且两算符可交换，就有 $e^T=e^{T_A}e^{T_B}$。对于适当可分离的参考态，波函数因此因子化，能量可相加。将指数代入本征方程并左乘 $e^{-T}$，得到
\begin{equation}
E=\langle\Phi_0|e^{-T}He^T|\Phi_0\rangle,\qquad
0=\langle\Phi_\mu|e^{-T}He^T|\Phi_0\rangle.
\end{equation}
CCSD 保留单、双激发投影方程；CCSD(T) 再对三激发作特定微扰修正。通常投影耦合簇不是变分上界法，近简并时单参考假设也可能失效。方法的声誉不能代替参考态适用性的检查。

\section{量子 Monte Carlo：对多体波函数采样}
对可求值的试探波函数 $\Psi_\theta(X)$，$X$ 表示全部电子坐标，能量可改写成
\begin{equation}
E(\theta)=\int p_\theta(X)E_L(X)\dd X,\quad
p_\theta=\frac{|\Psi_\theta|^2}{\int|\Psi_\theta|^2},\quad E_L=\frac{\hat H\Psi_\theta}{\Psi_\theta}.
\end{equation}
于是可以用第三章的采样思想计算积分，再优化参数 $\theta$。称为变分 Monte Carlo。真实本征态满足 $E_L=E$，局部能量方差为零；这是有价值的诊断。但有限试探族的最小能量通常仍高于精确值，统计估计还带误差。

另一条路是虚时演化，定义 $\tau=\ii t$ 后有 $-\partial_\tau\Psi=(H-E_T)\Psi$（原子单位）。本征态展开给出
\begin{equation}
\Psi(\tau)=\sum_n c_n e^{-(E_n-E_T)\tau}\Phi_n.
\end{equation}
只要初态与目标最低态重叠非零，归一化后高能成分相对衰减。动能在这一方程中像扩散项，势能像产生或消灭权重的项，形成扩散 Monte Carlo 的基础。费米波函数有正负号，不能直接作为概率密度。固定节点近似限制节点面为试探态的节点；对于局域 Hamilton 算符和理想投影，它给出固定节点能量上界，但节点误差不会靠增加样本自动消失。时间步、有限 walker 数以及非局域赝势处理还需独立控制。
